\chapter*{Phase 1--7 학습 지도}
\addcontentsline{toc}{chapter}{Phase 1--7 학습 지도}

\begin{factbox}
학습 순서는 ``텐서를 계산할 수 있음''에서 ``내부 알고리즘 가설을 인과적으로 시험할
수 있음''으로 진행한다. 뒤 Phase의 용어를 먼저 외우는 대신, 앞 Phase의 산출물을
다음 Phase의 측정 도구로 사용한다.
\end{factbox}

\begin{longtable}{@{}p{0.10\textwidth}p{0.32\textwidth}p{0.48\textwidth}@{}}
\caption{각 Phase의 핵심 질문과 통과 기준}\label{tab:phase-roadmap}\\
\toprule
Phase & 핵심 질문 & 통과 기준 \\
\midrule
\endfirsthead
\toprule
Phase & 핵심 질문 & 통과 기준 \\
\midrule
\endhead
1 & 벡터와 activation은 무엇인가? & affine map, dot product, norm, softmax,
embedding을 좌표와 기저의 언어로 설명하고 작은 예를 손으로 계산한다. \\
2 & 한 번의 순전파에서 무슨 일이 생기는가? & embedding부터 logits까지 모든
텐서의 shape와 index를 적고 attention 한 행과 MLP update를 직접 계산한다. \\
3 & 현대 LLM은 2017 Transformer와 무엇이 다른가? & causal decoder의
RMSNorm, RoPE, SwiGLU, GQA, KV cache와 prefill/decode를 구현 수준에서 설명한다. \\
4 & hidden state에서 무엇을 읽을 수 있는가? & probe, lens, geometry와 SAE가
측정하는 양을 구분하고 decodability와 causal use를 혼동하지 않는다. \\
5 & 구성요소는 어떻게 회로를 이루는가? & residual stream, QK/OV, head
composition, MLP feature를 이용해 induction 또는 IOI 회로의 경로를 추적한다. \\
6 & 회로 가설을 어떻게 반증하는가? & ablation, activation/path patching,
interchange intervention의 estimand와 통제 조건을 명시한다. \\
7 & 실제 행동의 메커니즘을 어디까지 아는가? & factual recall, ICL, reasoning,
refusal, hallucination 연구의 증거와 적용 범위를 비교하고 미해결 부분을 분리한다. \\
\bottomrule
\end{longtable}

\section*{두 개의 누적 프로젝트}

\subsection*{프로젝트 A: 순전파를 직접 소유하기}

Phase 1--3 동안 외부 Transformer 모듈을 호출하지 않는 최소 causal LM을 구현한다.
최종 코드는 적어도 다음 객체를 포함해야 한다.

\begin{lstlisting}[language=Python]
class RMSNorm
class RotaryEmbedding
class CausalSelfAttention
class SwiGLU
class TransformerBlock
class CausalLM
\end{lstlisting}

프로젝트 A의 통과 기준은 각 원소가 왜 그 값을 갖는지 설명할 수 있는가이다. 이를
검증하려면 각 연산에서 batch, sequence, head, feature 축을 출력하고, 작은 입력에서는
수작업 계산과 일치하는 unit test를 작성해야 한다.

\subsection*{프로젝트 B: 회로 가설을 시험하기}

프로젝트 B는 Phase 4--7 동안 공개된 소형 pretrained 모델에서 다음 검증 순서를 반복한다.

\begin{enumerate}
  \item 관찰할 behavior와 clean/corrupted prompt 분포를 고정한다.
  \item logits와 activation을 기록하고 descriptive hypothesis를 세운다.
  \item probe나 attribution으로 후보를 좁힌다.
  \item activation 또는 path intervention으로 causal effect를 측정한다.
  \item 발견한 회로의 faithfulness, completeness, minimality와 범위 밖 입력을 시험한다.
  \item 실패 사례에 따라 가설을 수정한다.
\end{enumerate}

\section*{권장 산출물}

각 Phase를 마칠 때 다음 네 파일을 남긴다.

\begin{itemize}
  \item 한 페이지의 정의와 수식 요약;
  \item 직접 유도한 계산 노트;
  \item seed와 환경을 기록한 최소 재현 코드;
  \item 이해한 내용, 반례, 열린 질문을 분리한 연구 기록.
\end{itemize}

이 책은 정답 모음이 아니라 이러한 산출물을 만들기 위한 기준 문서다. 최종 목표는
논문의 용어를 알아보는 것이 아니라, 논문의 intervention이 실제로 어느 tensor를
어떻게 바꾸며 그 결과가 어떤 인과 주장을 지지하는지 독립적으로 판단하는 것이다.
